\section{\blue{Proofs of the Regularizer and Optimality Properties}}
\label{app:optimality-properties}
\subsection{\blue{Proof of \Cref{lem:wtype-conj}}}
\label{app:wtype-conjugate-proof}
\begingroup\color{blue}
\begin{proof}
Let \(\psi(t):=\phi(\max\{|t|,1\})\). By \ref{W1}--\ref{W2}, \(\phi\) is convex and nonnegative on \([1,+\infty)\), with \(\phi(1)=0\); hence it is nondecreasing there. Thus \(s\mapsto\phi(\max\{s,1\})\) is convex and nondecreasing on \([0,+\infty)\), so \(\psi\) is convex and satisfies \(\psi\le\phi\). Conversely, any convex \(q\le\phi\) satisfies \(q(\pm1)\le0\), hence \(q\le0=\psi\) on \([-1,1]\), while \(q\le\phi=\psi\) on \(|t|\ge1\). Therefore \(\phi^{**}=\psi\), the largest lower semicontinuous convex minorant of \(\phi\)~\cite[Theorem~12.2]{rockafellar1970convex}. For \(|t|<1\), \ref{W1} gives \(\phi(t)>0=\psi(t)\), proving the contact-set formula.

Evaluating the definition of \(\phi^*\) at \(t=\pm1\) gives \(\phi^*(y)\ge|y|\). Since \(\phi(t_0)>0\) for any \(t_0>1\), convexity on \([1,+\infty)\) implies at least linear growth for large \(t\); evenness gives the same for negative \(t\). Hence \(\operatorname{dom}\phi^*\) contains an open interval around zero. The separable formulas follow from~\cite[Theorem~16.4]{rockafellar1970convex}.
\end{proof}
\endgroup

\subsection{Proof of \Cref{prop:zerogap}}
\label{app:zero-gap-proof}
\begin{proof}
    According to \cite[Theorem 2.158]{temporary-citekey-11039},
    $(\bx^{*}, \bB^{*}, \bLmbd^{*})$ are optimal solutions to the
    primal problem \eqref{Prop: block-wise regularized problem} and the dual problem
    \eqref{Prop: the dual problem} \blue{with zero duality gap} if and only if:
    \begin{align}
        \label{Eq: the sufficient and necessary condition}
        \begin{cases}
        (\bx^{*}, \bB^{*}) \in \argmin\limits_{\bx \in \RR^{q}, \bB \in \RR^{n\times d}}\calL(\blue{\bx, \bB}, \bLmbd^{*}), \\
        0 = h(\bB^{*}) + h^{*}(\bLmbd^{*}) - \langle \bLmbd^{*}, \bB^{*}\rangle.
        \end{cases}
    \end{align}
    The first condition in \eqref{Eq: the sufficient and necessary condition} implies
    that $\bm{0}$ belongs to the subdifferential of $\calL$ with respect to
    $(\bx, \bB)$ at $(\bx^{*}, \bB^{*}, \bLmbd^{*})$. Since $\calL$
    is continuously differentiable in $(\bx, \bB)$, we have
    \begin{align}
        \label{Eq: the sufficient and necessary condition for primal}
        \begin{cases}
        \bm{0} = \blue{\nabla_{\bx}F(\bx^{*}, \bB^{*})},\\
        \bm{0} = \blue{\nabla_{\bB}F(\bx^{*}, \bB^{*})} + \bLmbd^{*}.
        \end{cases}
    \end{align}
    For the second condition in \eqref{Eq: the sufficient and necessary condition},
    \blue{since $h^{**}\le h$, the Fenchel--Young inequality for $h^{**}$ yields
    $\langle\bLmbd^*,\bB^*\rangle\le h^{**}(\bB^*)+h^*(\bLmbd^*)\le h(\bB^*)+h^*(\bLmbd^*)=\langle\bLmbd^*,\bB^*\rangle$.
    Thus $h^{**}(\bB^*)+h^*(\bLmbd^*)=\langle\bLmbd^*,\bB^*\rangle$, which, by the subgradient inversion rule~\cite[Theorem~23.5]{rockafellar1970convex}, gives}
    \begin{align}
        \label{Eq: the sufficient and necessary condition for dual}\bB^{*}\in \partial h^{*}(\bLmbd^{*}).
    \end{align}
    Combining \eqref{Eq: the sufficient and necessary condition for primal} and \eqref{Eq: the sufficient and necessary condition for dual}, we obtain \blue{$\bm{0}\in\mathcal{K}(\bx^{*}, \bB^{*}, \bLmbd^{*})$ with $\mathcal{K}$ in \eqref{Eq: KKT mapping}}, completing the proof.
\end{proof}

\subsection{\blue{Proof of \Cref{prop:transfer}}}
\label{app:transfer-proof}
\begingroup\color{blue}
\begin{lemma}\label{lem:no-inner-stationary}
If \(\phi\) satisfies \ref{W1} and \ref{W3}, then \(0\notin\partial\phi(t)\) for every \(t\in(-1,1)\).
\end{lemma}
\begin{proof}
Fix \(t\in[0,1)\) and \(b\in(t,1)\). If \(t'\in[0,b)\) and \(v\in\hat\partial\phi(t')\), then \ref{W3} with \(s=t'\) and \(t'+h\), \(h\downarrow0\), gives \(v\le-c_b\). At \(t'=0\), evenness gives \(\phi(h)\le\phi(0)-c_b|h|\) for \(h<0\), hence also \(v\ge c_b\); thus \(\hat\partial\phi(0)=\emptyset\). By evenness, Fr\'echet subgradients at points of \((-b,0)\) are at least \(c_b\). Passing to limits, \(\partial\phi(t)\subset(-\infty,-c_b]\cup[c_b,\infty)\) for \(|t|<b\).
\end{proof}

\begin{proof}
    Since \(h^{*}\) is proper, lower semicontinuous, and convex, the subgradient inversion rule~\cite[Theorem~23.5]{rockafellar1970convex} gives
    \begin{align}
        \label{Eq: inversion rule of the conjugate}
        \tilde{\bB}\in\partial h^{*}(\tilde{\bLmbd})
        \iff \tilde{\bLmbd}\in\partial h^{**}(\tilde{\bB})
        \iff h^{**}(\tilde{\bB})+h^{*}(\tilde{\bLmbd})=\langle\tilde{\bLmbd},\tilde{\bB}\rangle.
    \end{align}
    The KKT conditions and \(\tilde{\bLmbd}=-\nabla_{\bB}F(\tilde{\bx},\tilde{\bB})\), \(\nabla_{\bx}F(\tilde{\bx},\tilde{\bB})=\bm{0}\), show that (a) holds if and only if \(\tilde{\bLmbd}\in\partial h(\tilde{\bB})\).

    (b)$\iff$(c). By Lemma~\ref{lem:wtype-conj}, the last equality in \eqref{Eq: inversion rule of the conjugate} is equivalent to \(h(\tilde{\bB})=h^{**}(\tilde{\bB})\), which holds if and only if \(|\tilde b_{ij}|\ge1\) for all \(i,j\).

    (b)$\Rightarrow$(a). By the definition of \(h^{*}\), (b) means that \(\tilde{\bB}\) minimizes \(h(\cdot)-\langle\tilde{\bLmbd},\cdot\rangle\) globally. Fermat's rule~\cite[Theorem~10.1]{Rockafellar1998VariationalA} and \cite[Exercise~8.8(c)]{Rockafellar1998VariationalA} give \(\tilde{\bLmbd}\in\partial h(\tilde{\bB})\).

    (a)$\Rightarrow$(c). By separability~\cite[Proposition~10.5]{Rockafellar1998VariationalA}, each entry satisfies \(\tilde\Lambda_{ij}\in\partial\phi(\tilde b_{ij})\cap\partial\phi^{**}(\tilde b_{ij})\). If \(|\tilde b_{ij}|<1\), Lemma~\ref{lem:wtype-conj} gives \(\partial\phi^{**}(\tilde b_{ij})=\{0\}\), so \(0\in\partial\phi(\tilde b_{ij})\), contradicting Lemma~\ref{lem:no-inner-stationary}. Hence \(|\tilde b_{ij}|\ge1\) for all \(i,j\).
\end{proof}
\endgroup

\noindent\blue{Our complexity bounds concern $\mathbb{E}[\mathrm{dist}^{2}(\bm{0},\mathcal K(\bx^{k},\bB^{k},\bLmbd^{k}))]$.} For $\bd^{k}\in\blue{\mathcal K}(\bx^{k},\bB^{k},\bLmbd^{k})$, let us denote $\bd^{k}=(\bd_{\bx}^{k},\bd_{\bB}^{k},\bd_{\bLmbd}^{k})$ and
\begin{align}
    \label{Eq: the components of the subgradient}
    \begin{cases}&\bd_{\bx}^{k}= \nabla_{\bx}F(\bx^{k}, \bB^{k}),\\&\bd_{\bB}^{k}=\nabla_{\bB}F(\bx^{k}, \bB^{k}) + \bLmbd^{k}, \\&\bd_{\bLmbd}^{k}= \bB^{k}- \calG_{\bLmbd}^{k},\end{cases}
\end{align}
where \blue{$\calG_{\bLmbd}^{1}:=\bB^1\in\partial h^*(\bLmbd^1)$ by proximal optimality in \eqref{eq:proximal-initialization} and subgradient inversion, and $\calG_{\bLmbd}^{k}:=2\bB^{k}-\bB^{k-1}+\tau(\bLmbd^{k-1}-\bLmbd^{k})\in\partial h^{*}(\bLmbd^{k})$ for $k\ge2$ by the optimality condition of the $\bLmbd$-update \eqref{Eq: the update of lambda}, implying}
\begin{align}
    \label{Eq: the stationarity error bound in proof}
    \mathrm{dist}^{2}(\bm{0},\blue{\mathcal K}(\bx^{k},\bB^{k},\bLmbd^{k}))\le\|\bd_{\bx}^{k}\|^{2}+\|\bd_{\bB}^{k}\|^{2}+\|\bd_{\bLmbd}^{k}\|^{2}.
\end{align}

\subsection{\blue{Proofs of \Cref{lem:conjugate-minorant,prop:original-certificate}}}
\label{app:original-certificate-proof}
\begingroup
\color{blue}
\paragraph{Conjugate minorant.}
Write $D=\mathcal{D}(\bx)$ and choose $V\in\partial h^*(\bLmbd)$ closest to $\bB$. Such a point exists because this set is nonempty and closed when $r_{\mathcal K}<\infty$. Set
\[
a=\nabla_{\bx}F(\bx,\bB),\qquad b=m(\bB-D)+\bLmbd,\qquad c=\bB-V.
\]
Then $r_{\mathcal K}^2=\|a\|^2+\|b\|^2+\|c\|^2$ and
$\nabla s_{\bLmbd}(\bx)=a+\nabla\mathcal D(\bx)^*b$, which proves~\eqref{eq:minorant-gradient}. The asserted Lipschitz constant follows from Assumption~\ref{Ass: basic regularity}(ii). Fenchel--Young gives $s_{\bLmbd}\le J_0$ and
$J_0(\bx)-s_{\bLmbd}(\bx)\ge h(D)-h^{**}(D)$.
Since $\bLmbd\in\partial h^{**}(V)$, Fenchel equality and $h^{**}\ge0$ yield
\[
\langle\bLmbd,D\rangle-h^*(\bLmbd)
=h^{**}(V)+\frac{\|\bLmbd\|^2}{m}
+\left\langle\bLmbd,c-\frac{b}{m}\right\rangle
\ge-\frac{\|mc-b\|^2}{4m}
\ge-\kappa r_{\mathcal K}^2.
\]
This proves~\eqref{eq:minorant-gap}.

\paragraph{Near-stationarity certificate.}
Let $s\le J_0$ have an $L$-Lipschitz gradient, where $L>0$, and suppose $J_0(\bx)-s(\bx)\le E$. Choose a global minimizer $\widehat{\bx}$ of
\[
H(\by):=J_0(\by)-\langle\nabla s(\bx),\by-\bx\rangle+L\|\by-\bx\|^2.
\]
The lower smoothness bound gives $H(\by)\ge s(\bx)+(L/2)\|\by-\bx\|^2$, so continuity of $J_0$ implies existence of this minimizer. Since $H(\widehat{\bx})\le H(\bx)=J_0(\bx)$,
\[
\|\widehat{\bx}-\bx\|\le\sqrt{2E/L}.
\]
Fermat's rule and the smooth sum rule give
$\nabla s(\bx)-2L(\widehat{\bx}-\bx)\in\partial J_0(\widehat{\bx})$, hence
\[
\operatorname{dist}(\bm0,\partial J_0(\widehat{\bx}))
\le\|\nabla s(\bx)\|+2\sqrt{2LE}.
\]
Applying this argument with $s=s_{\bLmbd}$ and $E=h(\mathcal{D}(\bx))+\kappa r_{\mathcal{K}}^2$, and using \Cref{lem:conjugate-minorant} proves part~(a), including $E=0$. For part~(b), $\bLmbd\in\dom h^*$ implies $\|\bLmbd\|\le\lambda\sqrt{nd}$.

\paragraph{Representative regularizers.}
For $u\in\RR$, $\log\cosh u\le|u|$ and $\operatorname{Hub}_{\delta}(u)\le|u|$. For the squared penalty, Assumption~\ref{Ass: basic regularity}(ii) gives
$\bigl||D_{ij}|-1\bigr|\le\max\{1,C_{\mathcal D}-1\}$.
Multiplying by $\bigl||D_{ij}|-1\bigr|$ and summing proves the bounds in \Cref{rem:certificate-table}.

\paragraph{Multiplier bound and the loss minorant.}
Define $T(U):=m(U-\operatorname{prox}_{h^{**}/m}(U))$. Firm nonexpansiveness of the convex proximal mapping implies that $T$ is $m$-Lipschitz. The residual identities give
$D=V+\bLmbd/m+c-b/m$, so $\bLmbd=T(D-c+b/m)$.
Since $h^{**}$ is separable, vanishes on the cube, and is nondecreasing away from it, each coordinate of $\operatorname{prox}_{h^{**}/m}(D)$ lies between $D_{ij}$ and its projection onto $[-1,1]$. Therefore
\[
\|\bLmbd\|\le\|T(D)\|+\|mc-b\|
\le m\|(|D|-\bm1)_+\|+\sqrt{1+m^2}\,r_{\mathcal K}.
\]
Using $\nabla f(\bx)=\nabla s_{\bLmbd}(\bx)-\nabla\mathcal D(\bx)^*\bLmbd$ and the preceding certificate argument with $s=f$, $L=L_f$, and $E=h(D)$ proves part~(i) of \Cref{rem:certificate-without-conjugate}. At an exact KKT point, $b=c=0$, giving $\bB=\operatorname{prox}_{h^{**}/m}(D)$ and $\bLmbd=T(D)$. For $p=1$, the scalar proximal formula gives $|\Lambda_{ij}|=\min\{m(|D_{ij}|-1)_+,\lambda\}$. On the cube, $T(D)=0$, $\bB=D$, and $\nabla f(\bx)=0$; also $h^*(0)=0$, so $s_{\bm0}=f$. This proves part~(ii).
\endgroup

\section{Auxiliary Results for the Complexity Analysis}
\label{app:complexity-auxiliary-results}

\subsection{\blue{Smoothness Constant of $F$}}
\label{app:smoothness-discussion}
\begingroup\color{blue}
Under Assumption~\ref{Ass: basic regularity}(ii), $F$ is $L_F$-smooth on $\RR^q\times\{\bB:\|\bB\|\le R_B\}$, where
\begin{equation}
\label{eq:derived-smoothness-constant}
L_F:=2L_f+\frac{\gamma}{n}\left(L_{\nabla\mathcal{D}}C_{\mathcal{D}B}+2G_{\mathcal{D}}^2+1\right).
\end{equation}
Here $C_{\mathcal{D}B}:=C_{\mathcal{D}}+R_B$ and
\[
R_B:=C_{\mathcal{D}}+\frac{n\lambda\sqrt{nd}}{\gamma}.
\]
This region contains all primal iterates under the prescribed initializations and step sizes, and $L_F$ is independent of $T$.

\begin{proof}
For a W-type regularizer with compact dual domain, $\bB^1\in\partial h^*(\bLmbd^1)$ and \eqref{Eq: the update of lambda} give $\|\bLmbd^k\|\le\lambda\sqrt{nd}$ for all $k\ge1$.
Since $L_F\ge\gamma/n$, the prescribed step sizes satisfy $0<\tau\gamma/n<1$.
Since $\phi^{**}$ is even and convex, its proximal mapping preserves sign and does not increase magnitude. Thus \eqref{eq:proximal-initialization} and Assumption~\ref{Ass: basic regularity}(ii) give $\|\bB^1\|\le\|\mathcal{D}(\bx^1)\|\le C_{\mathcal{D}}\le R_B$. If $\|\bB^k\|\le R_B$, then \eqref{Eq: the update of B} gives
\[
\|\bB^{k+1}\|\le\left(1-\frac{\tau\gamma}{n}\right)R_B
+\frac{\tau\gamma}{n}C_{\mathcal{D}}+\tau\lambda\sqrt{nd}=R_B,
\]
where the equality follows from the definition of $R_B$. By induction, $\|\bB^k\|\le R_B$ for all $k\ge1$.
Thus $\|\mathcal{D}(\bx)-\bB\|\le C_{\mathcal{D}B}$ throughout the stated region.

For two points $(\bx,\bB)$ and $(\bx',\bB')$ in this region, the gradient of $P(\bx,\bB)=\frac{\gamma}{2n}\|\mathcal{D}(\bx)-\bB\|^2$ satisfies
\begin{align*}
\|\nabla_{\bx}P(\bx,\bB)-\nabla_{\bx}P(\bx',\bB')\|
&\le\frac{\gamma}{n}\left[\left(L_{\nabla\mathcal{D}}C_{\mathcal{D}B}+G_{\mathcal{D}}^2\right)\|\bx-\bx'\|+G_{\mathcal{D}}\|\bB-\bB'\|\right],\\
\|\nabla_{\bB}P(\bx,\bB)-\nabla_{\bB}P(\bx',\bB')\|
&\le\frac{\gamma}{n}\left[G_{\mathcal{D}}\|\bx-\bx'\|+\|\bB-\bB'\|\right],
\end{align*}
by Assumption~\ref{Ass: basic regularity}(ii).
The coefficient matrix of these two bounds is positive semidefinite, so its spectral norm is at most its trace.
Hence $P$ is jointly smooth with constant $\frac{\gamma}{n}(L_{\nabla\mathcal{D}}C_{\mathcal{D}B}+G_{\mathcal{D}}^2+1)$.
Since $F=f+P$, \eqref{eq:derived-smoothness-constant} is a valid smoothness constant.
\end{proof}

\endgroup

\subsection{Proof of \Cref{Lemma: dual boundedness}}
\label{app:dual-difference-proof}
\begin{proof}
    Applying the $\bB$-update~\eqref{Eq: the update of B} at iterations $k+1$ and $k+2$ gives
    \begin{align*}
    \bLmbd^{k+1}-\bLmbd^k
    ={}&\tau^{-1}(\bB^{k+1}-\bB^{k+2})
    +\tau^{-1}(\bB^{k+1}-\bB^k)\\
    &+\nabla_{\bB}F(\bx^{k+1},\bB^k)
    -\nabla_{\bB}F(\bx^{k+2},\bB^{k+1}).
    \end{align*}
    By the $L_F$-smoothness of $F$ under \Cref{Ass: basic regularity}(ii),
    \begin{align*}
    \|\bLmbd^{k+1}-\bLmbd^k\|
    \leq{}&\tau^{-1}\|\bB^{k+2}-\bB^{k+1}\|
    +(\tau^{-1}+L_F)\|\bB^{k+1}-\bB^k\|\\
    &+L_F\|\bx^{k+2}-\bx^{k+1}\|.
    \end{align*}
    Squaring this inequality and using $(a+b+c)^2\leq3(a^2+b^2+c^2)$ proves
    \begin{align}
    \|\bLmbd^{k+1}-\bLmbd^k\|^2
    \leq{}&3\tau^{-2}\|\bB^{k+2}-\bB^{k+1}\|^2
    +3(\tau^{-1}+L_F)^2\|\bB^{k+1}-\bB^k\|^2\notag\\
    &+3L_F^2\|\bx^{k+2}-\bx^{k+1}\|^2,
    \label{Eq: dual boundedness in proof}
    \end{align}
    which is the claimed result.
\end{proof}

\subsection{Proof of \Cref{Lemma: One-iteration progress of B and Lambda}}
\label{app:primal-dual-progress-proof}
\begin{proof}
    \blue{By the descent lemma~\eqref{Eq: the descent lemma} for $F(\bx^{k+1},\cdot)$,
    \[
        F(\bx^{k+1},\bB^{k+1})-F(\bx^{k+1},\bB^{k})\le\langle\nabla_{\bB}F(\bx^{k+1},\bB^{k}),\bB^{k+1}-\bB^{k}\rangle+\frac{L_F}{2}\|\bB^{k+1}-\bB^{k}\|^{2}.
    \]
    The bilinear term satisfies
    \[
        \langle\bB^{k+1},\bLmbd^{k+1}\rangle-\langle\bB^{k},\bLmbd^{k}\rangle=\langle\bLmbd^{k},\bB^{k+1}-\bB^{k}\rangle+\langle\bB^{k+1},\bLmbd^{k+1}-\bLmbd^{k}\rangle,
    \]
    and the $\bB$-update gives
    \[
        \nabla_{\bB}F(\bx^{k+1},\bB^{k})+\bLmbd^{k}=-\tau^{-1}(\bB^{k+1}-\bB^{k}).
    \]
    For the conjugate term, $\calG_{\bLmbd}^{k}\in\partial h^{*}(\bLmbd^{k})$ defined after \eqref{Eq: the components of the subgradient} and the convexity of $h^{*}$ give
    \[
        -h^{*}(\bLmbd^{k+1})+h^{*}(\bLmbd^{k})\le-\langle\calG_{\bLmbd}^{k},\bLmbd^{k+1}-\bLmbd^{k}\rangle.
    \]
    Adding these three relations yields
    \begin{align*}
        &\calL(\bx^{k+1},\bB^{k+1},\bLmbd^{k+1})-\calL(\bx^{k+1},\bB^{k},\bLmbd^{k})\\
        &\le-\Bigl(\tau^{-1}-\frac{L_F}{2}\Bigr)\|\bB^{k+1}-\bB^{k}\|^{2}+\langle\bLmbd^{k+1}-\bLmbd^{k},\bB^{k+1}-\calG_{\bLmbd}^{k}\rangle.
    \end{align*}
    For $k\ge2$, the $\bB$-updates at iterations $k$ and $k-1$ and the expression of $\calG_{\bLmbd}^{k}$ give the identity below. For $k=1$, it follows from $\calG_{\bLmbd}^{1}=\bB^1$, $\bB^0=\bB^1$, and $\nabla_{\bB}F(\bx^1,\bB^1)+\bLmbd^1=\bm{0}$:
    \[
        \bB^{k+1}-\calG_{\bLmbd}^{k}=\tau\bigl(\nabla_{\bB}F(\bx^{k},\bB^{k-1})-\nabla_{\bB}F(\bx^{k+1},\bB^{k})\bigr),
    \]
    and hence}
    \begin{align}
        \label{Eq: the recursive relation of L on B and Lambda} & \calL(\bx^{k+1}, \bB^{k+1}, \bLmbd^{k+1}) - \calL(\bx^{k+1}, \bB^{k}, \bLmbd^{k}) \notag \\
        \leq & - \left(\tau^{-1}- \frac{L_{F}}{2}\right) \|\bB^{k+1}- \bB^{k}\|^{2}+ \Big\langle \bLmbd^{k+1}- \bLmbd^{k}, \tau \Big( \nabla_{\bB}F(\bx^{k}, \bB^{k-1}) - \nabla_{\bB}F(\bx^{k+1}, \bB^{k}) \Big) \Big\rangle.
    \end{align}
    For the inner product term in \eqref{Eq: the recursive relation of L on B and Lambda}, we obtain
    \begin{align}
        \label{Eq: the inner product in the recursive relation of L on B and Lambda} & \langle \bLmbd^{k+1}- \bLmbd^{k}, \tau (\nabla_{\bB}F(\bx^{k}, \bB^{k-1}) - \nabla_{\bB}F(\bx^{k+1}, \bB^{k})) \rangle \notag \\
        \leq & \, \frac{\delta\tau}{2}\|\bLmbd^{k+1}- \bLmbd^{k}\|^{2}+ \frac{\tau}{2\delta}\|\nabla_{\bB}F(\bx^{k+1}, \bB^{k}) - \nabla_{\bB}F(\bx^{k}, \blue{\bB^{k-1}}) \|^{2} \notag \\
        \leq & \, \frac{\delta\tau}{2}\|\bLmbd^{k+1}- \bLmbd^{k}\|^{2}+ \frac{\tau L_{F}^{2}}{2\delta}\left( \|\bx^{k+1}- \bx^{k}\|^{2}+ \|\bB^{k}- \bB^{k-1}\|^{2}\right),
    \end{align}
    where the first inequality follows from the fact that
    $\langle \bu, \bv \rangle \leq \frac{\delta}{2}\|\bu\|^{2}+ \frac{1}{2\delta}\|\bv\|^{2}$ for any $\delta >0$ and the second inequality follows from the $L_F$-smoothness of $F$. Then, \blue{setting $\delta=1/6$ and} substituting \eqref{Eq: the inner product in the recursive relation of L on B and Lambda}
    and \eqref{Eq: dual boundedness in proof} into \eqref{Eq: the recursive relation of L on B and Lambda}
    and rearranging the inequality yields \eqref{item: B and Lambda}.
\end{proof}

Before proceeding, throughout Appendices~\ref{app:stom-complexity} and~\ref{app:storm-complexity}, we adopt the notation, \blue{the probabilistic setting} and intermediate results from \Cref{Section: Convergence analysis}. 

\section{Complexity Analysis of DualHash-StoM (\texorpdfstring{\Cref{alg:our_algorithm_sgdm}}{Algorithm 1})}
\label{app:stom-complexity}

\blue{This appendix proves the complexity results for DualHash-StoM stated in \Cref{Section: complexity of StoM}.}
We first present a lemma on continuously differentiable functions with Lipschitz continuous gradients. Since the analysis also requires the relations among the extrapolated iterates, we next establish the key relationships among the iterates generated by \eqref{Eq: the update of x in SGDM} and \eqref{Eq: extrapolated points in SGDM}.

\begin{proposition}[\blue{Proposition 4.1 in \cite{InertialProximalAlternating}}]
    \label{Pp: relations of xyz} Let $\{\bx^{k}\}_{k \in \mathbb{N}}$,
    $\{\by^{k}\}_{k \in \mathbb{N}}$, and $\{\bz^{k}\}_{k \in \mathbb{N}}$ be sequences
    generated by \eqref{Eq: the update of x in SGDM} and \eqref{Eq: extrapolated points in SGDM}. For all $k \in \mathbb{N}$,
    we have:
    \begin{enumerate}[label=(\roman{*})]
        \item $\|\bx^{k}- \by^{k}\|^{2}= \alpha_{k}^{2}\|\bx^{k}- \bx^{k-1}\|^{2}$,

        \item $\|\bx^{k}- \bz^{k}\|^{2}= \beta_{k}^{2}\|\bx^{k}- \bx^{k-1}\|^{2}$,

        \item $\|\bx^{k+1}- \by^{k}\|^{2}\geq (1 - \alpha_{k})\|\bx^{k+1}- \bx^{k}
            \|^{2}+ \alpha_{k}(\alpha_{k}- 1)\|\bx^{k}- \bx^{k-1}\|^{2}$.
    \end{enumerate}
\end{proposition}

\begin{lemma}
    \label{Lemma: auxiliary lemma for f} Let $f(\bx): \mathbb{R}^{q}\rightarrow \mathbb{R}$
    be a continuously differentiable function with an $L$-Lipschitz continuous gradient.
    Then, for any $\bx,\, \by,\,\blue{\bg} \in \mathbb{R}^{q}$ and any $\bx^{+}\in \mathbb{R}
    ^{q}$ defined by
    \[
        \bx^{+}= \by - \eta \blue{\bg}, \hspace{5pt}\eta > 0,
    \]it holds that
    \begin{align}
        \label{Eq: auxiliary lemma for f}f(\bx^{+}) \leq f(\bx) + \langle \nabla f(\bx) - \blue{\bg}, \bx^{+}- \bx \rangle + \frac{L}{2}\|\bx^{+}- \bx\|^{2}+ \frac{1}{2\eta}\|\bx - \by\|^{2}- \frac{1}{2\eta}\|\bx^{+}- \by\|^{2}.
    \end{align}
\end{lemma}
\begin{proof}
    By the definition of $\bx^{+}$, we obtain that
    \[
        \bx^{+}= \argmin\limits_{\bu \in \mathbb{R}^{q}}\{ \langle \blue{\bg}, \bu
        - \by \rangle + \frac{1}{2\eta}\| \bu - \by\|^{2}\},
    \]
    and hence, in particular, by taking $\bu = \bx$, we obtain
    \begin{align*}
        0 \leq \langle \blue{\bg}, \bx - \bx^{+}\rangle + \frac{1}{2\eta}\| \bx - \by\|^{2}- \frac{1}{2\eta}\| \bx^{+}- \by\|^{2}.
    \end{align*}
    Invoking the descent lemma for a smooth function yields the result. 
\end{proof}

\subsection{Proof of \Cref{Lemma: One-iteration progress of x in SGDM}}
\label{app:stom-primal-progress-proof}

\begin{proof}
    Recalling the equivalent relation:
    \begin{align}
        \label{Eq: the update of x one progress in StoM}\mathcal{L}(\bx^{k+1}, \bm{B}^{k}, \bm{\Lambda}^{k}) - \mathcal{L}(\bx^{k}, \bm{B}^{k}, \bm{\Lambda}^{k}) = F(\bx^{k+1}, \bm{B}^{k}) - F(\bx^{k}, \bm{B}^{k}).
    \end{align}
    Then, applying \eqref{Eq: auxiliary lemma for f} \blue{to $F(\cdot, \bB^{k})$, which is $L_F$-smooth, with $\bx=\bx^{k}$, $\by=\by^{k}$, $\eta=\eta_{k}$, and $\bg=\bg_m^{k}$ (so that $\bx^{+}=\bx^{k+1}$ by \eqref{Eq: the update of x in SGDM}),}
    yields
    \begin{align}
        \label{Eq: the update of theta core}F(\bx^{k+1}, \bm{B}^{k}) \leq & F(\bx^{k}, \bm{B}^{k}) + \langle \nabla_{\bx}F(\bx^{k}, \bm{B}^{k}) - \blue{\bg_m^{k}}, \bx^{k+1}- \bx^{k}\rangle \notag                              \\
        +                                                                 & \frac{L_{F}}{2}\|\bx^{k+1}- \bx^{k}\|^{2}+ \frac{1}{2\eta_{k}}\|\bx^{k}- \bm{y}^{k}\|^{2}- \frac{1}{2\eta_{k}}\|\bx^{k+1}- \bm{y}^{k}\|^{2}.
    \end{align}
    For the inner product term in \eqref{Eq: the update of theta core}, using
    the facts that
    $\langle a, b \rangle \leq \frac{s}{2}\|a\|^{2}+ \frac{1}{2s}\|b\|^{2}$ for
    any $s > 0$ and $\| a - c \|^{2}\leq 2\|a - b \|^{2}+ 2\|b -c\|^{2}$ yields
    { \begin{align}\label{Eq: Inner product for theta}&\langle \nabla_{\bx}F(\bx^{k}, \bm{B}^{k}) - \blue{\bg_m^{k}}, \bx^{k+1}- \bx^{k}\rangle \notag \\ \leq&\,\frac{s}{2}\|\bx^{k+1}- \bx^{k}\|^{2}+ \frac{1}{s}\|\nabla_{\bx}F(\bx^{k}, \bm{B}^{k}) - \nabla_{\bx}F(\bm{z}^{k}, \bm{B}^{k}) \|^{2}+ \frac{1}{s}\|\nabla_{\bx}F(\bm{z}^{k}, \bm{B}^{k}) - \blue{\bg_m^{k}}\|^{2}\notag \\ \leq&\, \frac{s}{2}\|\bx^{k+1}- \bx^{k}\|^{2}+ \frac{L^{2}_{F}}{s}\|\bx^{k}- \bm{z}^{k}\|^{2}+ \frac{1}{s}\|\nabla_{\bx}F(\bm{z}^{k}, \bm{B}^{k}) - \blue{\bg_m^{k}}\|^{2}.\end{align}}
    Setting $s = L_{F}$ in \eqref{Eq: Inner product for theta} and substituting into
    \eqref{Eq: the update of theta core}, then applying \cref{Pp: relations of xyz},
    we obtain
    \begin{align}
        F(\bx^{k+1}, \bm{B}^{k}) & \leq F(\bx^{k}, \bm{B}^{k}) - (\frac{1 - \alpha_{k}}{\eta_{k}}- 2L_{F})\frac{\|\bx^{k+1}- \bx^{k}\|^{2}}{2}\notag                                                                                           \\
                                 & +\,(\frac{\alpha_{k}}{\eta_{k}}+ 2L_{F}\beta_{k}^{2})\frac{\|\bx^{k}- \bx^{k-1}\|^{2}}{2}+ \frac{1}{L_{F}}\|\nabla_{\bx}F(\bm{z}^{k}, \bm{B}^{k}) - \blue{\bg_m^{k}}\|^{2}. \label{Eq: the update of x of L in StoM}
    \end{align}
    Substituting \eqref{Eq: the update of x of L in StoM} into \eqref{Eq: the update of x one progress in StoM}
    yields \eqref{Eq: One-iteration progress of x in SGDM} and completes the proof.
    
    % \begin{align}
    %     \label{item: x}
    %     &\mathcal{L}(\bx^{k+1}, \bm{B}^k, \bm{\Lambda}^k) -\mathcal{L}(\bx^k, \bm{B}^k, \bm{\Lambda}^k)  \notag \\
    %     \leq &\,-(\frac{1 - \alpha_k}{\eta_k}- 2L_F) \frac{ \|\bx^{k+1} - \bx^k \|^2}{2}
    %     + (\frac{\alpha_k}{\eta_k}+ 2L_F\beta_k^2) \frac{\|\bx^k - \bx^{k-1} \|^2}{2} \notag \\
    %     \quad&+ \frac{1}{L_F} \|\nabla_{\bx} F(\bm{z}^k, \bm{B}^k) - \bG^k \|^2.
    % \end{align} 
\end{proof}
\subsection{Proof of \Cref{Lemma: approximate sufficient descent property in SGDM}}
\label{app:stom-descent-proof}
\begin{proof}
    Substituting one-iteration progress \eqref{Eq: One-iteration progress of x in SGDM}
    and \eqref{item: B and Lambda} into \eqref{Eq: One-iteration progress decomposition}
    , we obtain
    \begin{align}
        \label{Eq: One-iteration progress in sgdm} & \calL (\bx^{k+1}, \bB^{k+1}, \bLmbd^{k+1}) - \calL (\bx^{k}, \bB^{k}, \bLmbd^{k}) \notag                                                                                                   \\
        \leq \,                                    & K_{1}\Delta_{\bB}^{k+2}- K_{2}\Delta_{\bB}^{k+1}+ K_{3}\Delta_{\bB}^{k}+ K_{4}\Delta_{\bx}^{k+2}- (K_{5}- K_{3})\Delta_{\bx}^{k+1}+ K_{6}\Delta_{\bx}^{k}+ \frac{1}{L_{F}}\|\be^{k}\|^{2}.
    \end{align}
    Applying the definition of $\Psi_m$ from \eqref{Eq: the Lyapunov function in SGDM}
    to \eqref{Eq: One-iteration progress in sgdm}, rearranging, and taking the
    expectation on both sides, we obtain
    {\small
    \begin{align}
        \label{the descent property of Lyapunov func in SGDM first} & \EE\left[\Psi_m(\mathcal{Q}^{k+1})\right] - \EE\left[\Psi_m(\mathcal{Q}^{k})\right]\notag                                                                               \\
        \leq                                                        & - (C_{1}- K_{1}) \EE[\Delta_{\bm{B}}^{k+2}] -(K_{2}- C_{1}- C_{2}) \EE[\Delta_{\bm{B}}^{k+1}] - (C_{2}- K_{3}) \EE[\Delta_{\bm{B}}^{k}] \notag                     \\
                                                                    & - (C_{3}-K_{4}) \EE[\Delta_{\bx}^{k+2}]- (K_{5}-K_{3}-C_{3}-C_{4}) \EE[\Delta_{\bx}^{k+1}] - (C_{4}- K_{6}) \EE[\Delta_{\bx}^{k}] + \frac{\sigma^{2}}{b_{k}L_{F}}.
    \end{align}}Then, \blue{by \Cref{lem:lyapunov-coef-positivity},} substituting the values of coefficients in \eqref{Eq: the coefficients of the Lyapunov function in SGDM}
    and parameter conditions \eqref{Eq: the parameters conditions of dualhash-sgdm in suppl}
    into \eqref{the descent property of Lyapunov func in SGDM first}, we derive
    \begin{align*}
             & \EE\left[\Psi_m(\mathcal{Q}^{k+1})\right] - \EE\left[\Psi_m(\mathcal{Q}^{k})\right] \nonumber                                                                                                                                                        \\
        \leq & \, -C_{\bB}\left( \EE\left[\Delta_{\bB}^{k+1}\right] + \EE\left[\Delta_{\bB}^{k}\right] \right) -C_{\bx}\left( \EE\left[\Delta_{\bm{x}}^{k+1}\right]+ \EE\left[\Delta_{\bm{x}}^{k}\right] \right) + \frac{\sigma^{2}}{b_{k}L_{F}}.
    \end{align*}which is exactly \eqref{Eq: the descent coefficients of potential function in SGDM},
    thereby completing the proof. 
\end{proof}
The technical proof that the coefficients of $\Psi_m$ in \eqref{Eq: the coefficients of the Lyapunov function in SGDM} and the descent constants $C_{\bB}$, $C_{\bx}$ are indeed positive is given as follows.
\begin{lemma}[Positivity of Lyapunov related coefficients]
    \label{lem:lyapunov-coef-positivity} If parameter conditions \eqref{Eq: the parameters conditions of dualhash-sgdm in suppl}
    hold, then the following inequalities are satisfied:
    \begin{align}
        % \label{Eq: the positive coefficients}
        \begin{cases}K_{5}- K_{3}- C_{3}- C_{4}\geq \nu C_{4}, \\ C_{4}- K_{6}\geq \nu C_{4},\\ C_{4}> 0,\\ C_{\bB}> 0 ,\\ C_{2},\, C_{1}> 0, \\ C_{\bx}> 0.\end{cases} \notag
    \end{align}
\end{lemma}
\begin{proof}
    For $K_{5}- K_{3}- C_{3}- C_{4}$ and $C_{4}- K_{6}$ in \eqref{the descent property of Lyapunov func in SGDM first},
    we obtain by straightforward computations by choosing some $\nu > 0$\blue{, where we recall that $K_{5}= \frac{1 - \alpha_{k}}{\eta_{k}}- 2L_{F}$, $C_{3}= K_{4}$, and $\bar{L}_{F}= 2L_{F}+ K_{3}+ K_{4}$}:
    {\begin{align}
        K_{5}- K_{3}- C_{3}- C_{4}
        &= \frac{1 - \alpha_{k}}{\eta_{k}}- \bar{L}_{F}- (1 + \nu)C_{4}+ \nu C_{4}\notag\\
        &= (1 - \alpha_{k})\left(\frac{1}{\eta_{k}}- \frac{\bar{L}_{F}+ (1 + \nu)C_{4}}{1 - \alpha_{k}}\right) + \nu C_{4}\notag \\
        &\geq \nu C_{4}. \notag % \label{Eq: the scaling of the delta1}
    \end{align}}
    where the last inequality follows from $\alpha_{k}= \alpha$ and
    $\frac{1}{\eta_{k}}= \frac{\bar{L}_{F}+ (1 + \nu)C_{4}}{1 - \alpha}$. Then, \blue{using $\alpha_{k}= \alpha$ and $\beta_{k}= \beta$ from \eqref{Eq: the parameters conditions of dualhash-sgdm in suppl},}
    one has
    \begin{align}
        C_{4}- K_{6} & = C_{4}- \frac{\alpha}{\eta_{k}}- 2L_{F}\beta_{k}^{2}\notag                                                                                                                                   \\
                     & = \frac{1}{1 - \alpha}\left( (1 - \alpha)C_{4}- \alpha\bar{L}_{F}- (1 + \nu)C_{4}\alpha - 2L_{F}\beta^{2}(1 - \alpha) \right)\notag                                                           \\
                     & \blue{=} \frac{1}{1 - \alpha}\left( (1 - 2\alpha - \nu)C_{4}- \alpha \left(\bar{L}_{F}+ \frac{2L_{F}\beta^{2}}{\alpha}\right) \right) + \nu C_{4}+ \frac{2L_{F}\beta^{2}\alpha}{1 - \alpha}\notag \\
                     & \geq \nu C_{4},\notag % \label{Eq: the scaling of the delta2}
    \end{align}where the last inequality follows from $\alpha \in (0, \frac{1 - \nu}{2})$
    and
    $C_{4}= \frac{\bar{L}_{F}+ \frac{2L_F \beta^2}{\alpha}}{1 - 2\alpha - \nu}> 0$.
    We rearrange $\frac{1}{\eta_{k}}$ using the definition of $C_{4}$ and obtain
    $\eta_{k}= \frac{\eta}{L_{F}}$ where $\eta = \frac{(1 - 2\alpha - \nu)/2}{\frac{\bar{L}_{F}}{L_{F}}
    + \frac{(1 + \nu)\beta^{2}}{(1 - \alpha)\alpha}}$.

    Next, we verify that the coefficients of $\Psi_m$ satisfying \eqref{Eq: the coefficients of the Lyapunov function in SGDM}
    and $\tau$ satisfying \eqref{Eq: the parameters conditions of dualhash-sgdm in suppl}
    ensure that $C_{\bB},\, C_{\bx}> 0$. Let $t:=\tau L_F$. \blue{By the definitions of $K_1,K_2,K_3$,}
    \[
        \blue{C_{\bB}=\frac{K_2-K_3-K_1}{2}=\frac{L_F}{2t}\left(1-2t-\frac{13}{2}t^2\right).}
    \]
    \blue{The quadratic $1 - 2t - \frac{13}{2}t^{2}$ is positive if and only if $0 < t < \frac{\sqrt{30}-2}{13}$. Hence selecting
    $\tau \in \left(0, \frac{\sqrt{30}-2}{13L_{F}}\right)$ ensures $C_{\bB}> 0$. Consequently, $C_1=K_1=\frac{1}{2\tau}>0$ and $C_2=\frac{1}{2}(K_2-K_1+K_3)=C_{\bB}+K_3>0$, since $K_3=6\tau L_F^2>0$.}
    \blue{Finally, $C_3=K_4>0$ and $C_{\bx}=\nu C_4>0$. Thus all the stated coefficients are positive under \eqref{Eq: the parameters conditions of dualhash-sgdm in suppl}.}
\end{proof}
\subsection{\blue{Lower Bound of the StoM Potential}}
\label{app:stom-lower-bound}
\begingroup\color{blue}
We first record two estimates that are used here and in \Cref{app:storm-lower-bound}.
Let $t:=\tau L_F$. Both parameter settings \eqref{Eq: the parameters conditions of dualhash-sgdm in suppl} and \eqref{Eq: the parameters conditions in dualhash-storm}--\eqref{Eq: the constants for storm} give $t<\frac{\sqrt{30}-2}{13}<0.27$. Hence, by \eqref{Eq: the coefficients of the Lyapunov function in SGDM},
\[
C_1=K_1=\frac{1}{2\tau},\qquad C_3=K_4=\frac{\tau L_F^2}{2}.
\]
By Appendix~\ref{app:smoothness-discussion}, $\|\bB^k\|\le R_B$ and $\|\bLmbd^k\|\le\lambda\sqrt{nd}$ for all $k\ge1$, and $0\le h^*\le nd\sup_{y\in\operatorname{dom}\phi^*}\phi^*(y)$ on $\operatorname{dom}h^*$ by \Cref{lem:wtype-conj} and the condition on $\phi^*$ in \Cref{Section: W-regularizer}. Therefore $\langle\bB^k,\bLmbd^k\rangle\ge-\lambda\sqrt{nd}\,R_B$ and $\Delta_{\bB}^{k+1}=\frac12\|\bB^{k+1}-\bB^k\|^2\le2R_B^2$, so for both algorithms and every $k\ge1$,
\begin{equation}
\label{eq:lb-common}
\calL(\bx^k,\bB^k,\bLmbd^k)-C_1\Delta_{\bB}^{k+1}
\ge F(\bx^k,\bB^k)-c_1',
\qquad
c_1':=\lambda\sqrt{nd}\,R_B+nd\sup_{y\in\operatorname{dom}\phi^*}\phi^*(y)+\frac{R_B^2}{\tau}.
\end{equation}
It remains to control $C_3\Delta_{\bx}^{k+1}=\frac{\tau L_F^2}{4}\|\bx^{k+1}-\bx^k\|^2$ by $F-\underline{\ell}$. Since $f\ge\underline{\ell}$ by \Cref{Ass: basic regularity}(i) and $P\ge0$, we have $F\ge\underline{\ell}$. Let $\Omega:=\RR^q\times\{\bB:\|\bB\|\le R_B\}$. For $(\bx,\bB)\in\Omega$, the point $(\bx,\bB)-L_F^{-1}\nabla F(\bx,\bB)$ has $\bB$-component $(1-\frac{\gamma}{nL_F})\bB+\frac{\gamma}{nL_F}\mathcal{D}(\bx)$, a convex combination of $\bB$ and $\mathcal{D}(\bx)$ since $L_F\ge\gamma/n$; it therefore lies in the convex set $\Omega$ because $R_B\ge C_{\mathcal{D}}$. The descent lemma on $\Omega$ and $F\ge\underline{\ell}$ give
\begin{equation}
\label{eq:lb-gradient}
\|\nabla F(\bx,\bB)\|^2\le2L_F\bigl(F(\bx,\bB)-\underline{\ell}\bigr),\qquad(\bx,\bB)\in\Omega .
\end{equation}

\begin{lemma}[Lower bound of $\Psi_m$]
\label{Lemma: lower bound of Psi_m}
Under Assumptions~\ref{Ass: basic regularity} and~\ref{Ass: main assumptions} and the parameter conditions \eqref{Eq: the parameters conditions of dualhash-sgdm in suppl}, for every $k\ge1$,
\[
\EE[\Psi_m(\calQ^k)]\ge\Psi_m^*:=\underline{\ell}-c_1,
\qquad
c_1:=c_1'+\frac34\tau\eta^2\sigma^2,
\]
where $c_1'$ is defined in \eqref{eq:lb-common}; $c_1$ is independent of $T$.
\end{lemma}

\begin{proof}
Write $\bm{p}_k:=\bx^k-\bx^{k-1}$. By \eqref{Eq: the update of x in SGDM} and the definition of $\be^k$,
\[
\bx^{k+1}-\bx^k=\alpha\bm{p}_k-\frac{\eta}{L_F}\bigl(\nabla_{\bx}F(\bz^k,\bB^k)+\be^k\bigr).
\]
Since $\|\bB^k\|\le R_B$, the $L_F$-smoothness of $F$ and \eqref{eq:lb-gradient} give
\[
\|\nabla_{\bx}F(\bz^k,\bB^k)\|^2\le2\|\nabla_{\bx}F(\bx^k,\bB^k)\|^2+2L_F^2\beta^2\|\bm{p}_k\|^2
\le4L_F\bigl(F(\bx^k,\bB^k)-\underline{\ell}\bigr)+2L_F^2\beta^2\|\bm{p}_k\|^2 .
\]
Combining these two displays with $\EE\|\be^k\|^2\le\sigma^2$ from \eqref{eq: the gradient error in SGDM} yields
\[
C_3\EE[\Delta_{\bx}^{k+1}]
\le3t\eta^2\,\EE\bigl[F(\bx^k,\bB^k)-\underline{\ell}\bigr]
+\frac{tL_F}{4}\left(3\alpha^2+6\eta^2\beta^2\right)\EE\|\bm{p}_k\|^2
+\frac34\tau\eta^2\sigma^2 .
\]
Under \eqref{Eq: the parameters conditions of dualhash-sgdm in suppl}, $\bar L_F\ge2L_F$ gives $\eta<\frac14$, and $0<1-2\alpha-\nu<1$ gives $C_4\ge\bar L_F\ge2L_F$. With $\alpha<\frac12$ and $\beta<1$,
\[
3t\eta^2<1,
\qquad
\frac{tL_F}{4}\left(3\alpha^2+6\eta^2\beta^2\right)<\frac{9tL_F}{32}<L_F\le\frac{C_4}{2}.
\]
Hence the term $C_4\Delta_{\bx}^k=\frac{C_4}{2}\|\bm{p}_k\|^2$ dominates the $\|\bm{p}_k\|^2$-term. Using \eqref{eq:lb-common} and $C_2\Delta_{\bB}^k\ge0$,
\[
\EE[\Psi_m(\calQ^k)]
\ge\underline{\ell}+\left(1-3t\eta^2\right)\EE\bigl[F(\bx^k,\bB^k)-\underline{\ell}\bigr]-c_1'-\frac34\tau\eta^2\sigma^2
\ge\underline{\ell}-c_1 .
\]
\end{proof}
\endgroup

\subsection{Proof of \Cref{lemma: convergence rate sgdm}}
\label{app:stom-rate-proof}
We then bound the three terms on the right-hand side of \eqref{Eq: the stationarity error bound in proof}.

\begin{lemma}[Stationarity error bound]
    \label{Lemma: Stationarity error bound in SGDM} Under Assumptions~\ref{Ass: basic regularity} and~\ref{Ass: main assumptions},
    for any $k \geq 1$, it holds that
    \begin{align}\label{Eq: the stationarity error bound in SGDM Lemma}
        \EE [\mathrm{dist}^{2}(\bm{0}, \mathcal K(\bx^{k}, \bB^{k}, \bm{\Lambda}^{k}))] \leq S_{1}\EE[ \Delta_{\bB}^{k+1}] + S_{2}\EE[ \Delta_{\bB}^{k}] + S_{3}^{k}\EE[ \Delta_{\bm{x}}^{k+1}] + S_{4}^{k}\EE [\Delta_{\bm{x}}^{k}] + \blue{\frac{3\sigma^{2}}{b_{k}}},
    \end{align}
    where the coefficients are given by
    \begin{align}
        \label{Eq: the coefficients of stationarity error bound in SGDM}\begin{cases}S_{1}&= \blue{4(\tau^{-2}+ 3)},\\ S_{2}&= 4(1 + 3 (1 + \tau L_{F})^{2}),\\ S_{3}^{k}&= 4(3\eta_{k}^{-2}+ L_{F}^{2}+ 3 (\tau L_{F})^{2}),\\ S_{4}^{k}&= 6 ( L^2_{F}\beta_{k}^{2}+ \blue{2} \eta_{k}^{-2}\alpha_{k}^{2}).\end{cases}
    \end{align}
\end{lemma}
% \begin{proof}
% For estimating \(\bd_{\bx}^k\), using the optimality condition of \(\bx^k\), we obtain
% \begin{align}
%     &\|\bd_{\bx}^k\|^2 \notag \\
%     =&\,\|\nabla_{\bx} F(\bx^k, \bm{B}^k) \|^2  \notag\\
%     =&\,\eta_k^{-2} \left\| \eta_k \nabla_{\bx} F(\bx^k, \bm{B}^k) + 0\right\|^2 \notag \\
%     =&\, \eta_k^{-2}  \left\|\eta_k\left( \nabla_{\bx} F(\bx^k, \bm{B}^k) - \bG^k \right) + \bm{y}^k - \bx^{k+1} \right\|^2 \notag \\
%     =&\,\eta_k^{-2} \bigg\|\eta_k\Big(
%     \nabla_{\bx} F(\bx^k, \bm{B}^k) - \nabla_{\bx} F(\bm{z}^k, \bm{B}^k) \notag \\
%     &\qquad + \nabla_{\bx} F(\bm{z}^k, \bm{B}^k)- \bG^k \Big)
%     + \bm{y}^k - \bx^k + \bx^k - \bx^{k+1} \bigg\|^2
%     \label{eq:ineq1} \\
%     \leq&\, 3 \|\nabla_{\bx} F(\bx^k, \bm{B}^k) - \nabla_{\bx} F(\bm{z}^k, \bm{B}^k)\|^2
%     + 3\eta_k^{-2}\|\bx^{k+1} - \bx^k + \bx^k - \bm{y}^k \|^2 \notag \\
%     \quad& + 3 \|\be^k\|^2
%     \label{eq:ineq2} \\
%     \leq&\, 3L_F^2 \|\bx^k - \bm{z}^k \|^2
%     + 6\eta_k^{-2} \left(\|\bx^{k+1} - \bx^k\|^2 + \|\bx^k - \bm{y}^k \|^2\right) \notag \\
%     \quad& + 3 \|\be^k\|^2,
% \end{align} where \eqref{eq:ineq1} is deduced by using the fact that \((a + b+ c)^2 \leq 3(a^2 + b^2 +c^2)\) and \eqref{eq:ineq2} is deduced by the smoothness of \( F(\bx, \bm{B})\) and \((a + b)^2 \leq 2(a^2 + b^2)\).
% Then using \cref{Pp: relations of xyz} concludes
%     \begin{align}
%     \label{Eq: dx}
%     \| \bd_{\bx}^k\| \leq  6( L_F \beta_k^2 + 2 \eta_k^{-2} \alpha_k^2) \Delta_{\bx}^{k} + 12 \eta_k^{-2} \Delta_{\bx}^{k+1} + 3 \| \be^k\|^2
%     \end{align}
%     For estimating  \(\bd^k_{\bB}\), using the optimality condition of \(\bB^k\) and the smoothness of \(F\), we obtain
% \begin{align}
%     \|\bd_{\bB}^k\|^2&= \| \nabla_{\bB} F(\bx^{k+1}, \bB^k)+  \bm{\Lambda}^k + \nabla_{\bB} F(\bx^k, \bB^k) - \nabla_{\bB} F(\bx^{k+1}, \bB^k) \|^2 \notag \\
%                     &= 4 \tau^{-2} \Delta^{k+1}_{\bB}  + 4L_F^2\Delta_{\bx}^{k+1} \label{Eq: dB},
% \end{align} where the last inequality follows from the fact \((a + b)^2 \leq 2(a^2 + b^2)\).

% For estimating \( \|\bd^k_{\bm{\Lambda}}\|^2\),  using the optimality condition of \( \bm{\Lambda}^k \) in \eqref{Eq: the optimality condition of lbda},  we obtain
% \begin{align*}
%         \|\bd^k_{\bm{\Lambda}}\|^2 &\leq 2\tau^{2} \|\bm{\Lambda}^k - \bm{\Lambda}^{k-1} \|^2+ 2\|\bB^{k} - \bB^{k-1}\|^2.
%     \end{align*} Then, substituting \eqref{Eq: dual boundedness} in \cref{Lemma: dual boundedness} into the above inequality yield that
% \begin{align}
% \label{Eq: dLambda}
% \|\bd^k_{\bm{\Lambda}} \|^2 \leq 12 \Delta_{\bB}^{k+1} + 4(1 + 3 (1 + \tau L_F )^2)\Delta_{\bB}^k + 12(\tau L_F)^2\Delta_{\bx}^{k+1} .
% \end{align}
% Then, substituting \eqref{Eq: dx}, \eqref{Eq: dB} and \eqref{Eq: dLambda} into \eqref{Eq: the stationarity error bound in proof} yields \eqref{Eq: the stationarity error bound in SGDM Lemma} and completes the proof. 
% \end{proof}
\begin{proof}
    For estimating $\| \bd_{\bx}^{k}\|^{2}$, using the optimality condition of
    $\bx^{k}$, we obtain
    {\small
    \begin{align}
             & \|\bd_{\bx}^{k}\|^{2}\notag                                                                                                                                                                                                                   \\
        =    & \,\|\nabla_{\bx}F(\bx^{k}, \bm{B}^{k}) \|^{2}\notag                                                                                                                                                                                           \\
        =    & \,\eta_{k}^{-2}\left\| \eta_{k}\nabla_{\bx}F(\bx^{k}, \bm{B}^{k}) + 0\right\|^{2}\notag                                                                                                                                                       \\
        =    & \, \eta_{k}^{-2}\left\|\eta_{k}\left( \nabla_{\bx}F(\bx^{k}, \bm{B}^{k}) - \blue{\bg_m^{k}}\right) + \bm{y}^{k}- \bx^{k+1}\right\|^{2}\notag                                                                                                           \\
        =    & \,\eta_{k}^{-2}\bigg\|\eta_{k}\Big( \nabla_{\bx}F(\bx^{k}, \bm{B}^{k}) - \nabla_{\bx}F(\bm{z}^{k}, \bm{B}^{k}) + \nabla_{\bx}F(\bm{z}^{k}, \bm{B}^{k})- \blue{\bg_m^{k}}\Big) + \bm{y}^{k}- \bx^{k}+ \bx^{k}- \bx^{k+1}\bigg\|^{2}\label{eq:d_x_ineq1} \\
        \leq & \, 3 \|\nabla_{\bx}F(\bx^{k}, \bm{B}^{k}) - \nabla_{\bx}F(\bm{z}^{k}, \bm{B}^{k})\|^{2}+ 3\eta_{k}^{-2}\|\bx^{k+1}- \bx^{k}+ \bx^{k}- \bm{y}^{k}\|^{2}+ 3 \|\be^{k}\|^{2}\label{eq:d_x_ineq2}                                                 \\
        \leq & \, 3L_{F}^{2}\|\bx^{k}- \bm{z}^{k}\|^{2}+ 6\eta_{k}^{-2}\left(\|\bx^{k+1}- \bx^{k}\|^{2}+ \|\bx^{k}- \bm{y}^{k}\|^{2}\right) + 3 \|\be^{k}\|^{2},\label{eq:d_x_ineq3}                                                                      
    \end{align}}where \eqref{eq:d_x_ineq2} is deduced by using the fact that
    $(a + b+ c)^{2}\leq 3(a^{2}+ b^{2}+c^{2})$ and \eqref{eq:d_x_ineq3} is deduced
    \blue{from the $L_F$-smoothness of $F(\cdot, \bB^{k})$} and $(a + b)^{2}\leq 2(a^{2}+ b^{2})$.
    Then applying \Cref{Pp: relations of xyz} to \eqref{eq:d_x_ineq3} concludes
    \begin{align}
        \label{Eq: dx}\| \bd_{\bx}^{k}\|^{2}\leq 6( L^2_{F}\beta_{k}^{2}+ 2 \eta_{k}^{-2}\alpha_{k}^{2}) \Delta_{\bx}^{k}+ 12 \eta_{k}^{-2}\Delta_{\bx}^{k+1}+ 3 \| \be^{k}\|^{2}.
    \end{align}
    For estimating $\|\bd^{k}_{\bB}\|^{2}$, using the $\bB$-update
    \eqref{Eq: the update of B}, \blue{the $L_F$-smoothness of $F$, and the inequality $(a + b)^{2}\leq 2(a^{2}+ b^{2})$}, we obtain
    \begin{align}
        \|\bd_{\bB}^{k}\|^{2} & = \| \nabla_{\bB}F(\bx^{k+1}, \bB^{k})+ \bm{\Lambda}^{k}+ \nabla_{\bB}F(\bx^{k}, \bB^{k}) - \nabla_{\bB}F(\bx^{k+1}, \bB^{k}) \|^{2}\notag \\
                              & \leq4 \tau^{-2}\Delta^{k+1}_{\bB}+ 4L_{F}^{2}\Delta_{\bx}^{k+1}.\label{Eq: dB}
    \end{align}

    \blue{For $k=1$, the initialization gives $\bd_{\bLmbd}^{1}=\bm{0}$, so \eqref{Eq: dLambda} holds. For $k\ge2$, using the $\bm{\Lambda}$-update \eqref{Eq: the update of lambda}, we obtain}
\begin{align}
    \|\bd^{k}_{\bm{\Lambda}}\|^{2} & \leq 2\tau^{2}\|\bm{\Lambda}^{k}- \bm{\Lambda}^{k-1}\|^{2}+ 2\|\bB^{k}- \bB^{k-1}\|^{2}. \notag
\end{align}

    Then, \blue{bounding $\|\bm{\Lambda}^{k}- \bm{\Lambda}^{k-1}\|^{2}$ by \Cref{Lemma: dual boundedness} (with $k$ replaced by $k-1$)}, we obtain
    \begin{align}
        \label{Eq: dLambda}\|\bd^{k}_{\bm{\Lambda}}\|^{2}\leq 12 \Delta_{\bB}^{k+1}+ 4(1 + 3 (1 + \tau L_{F})^{2})\Delta_{\bB}^{k}+ 12(\tau L_{F})^{2}\Delta_{\bx}^{k+1}.
    \end{align}

    Finally, \blue{we substitute} \eqref{Eq: dx}, \eqref{Eq: dB}, and \eqref{Eq: dLambda}
    into \eqref{Eq: the stationarity error bound in proof} \blue{and take} the full expectation
    on both sides, \blue{which leads to} \eqref{Eq: the stationarity error bound in SGDM Lemma},
    completing the proof. 
\end{proof}

% We now establish the convergence rate of \Cref{alg:our_algorithm_sgdm} as follows.
% We set $T$ as the maximum number of iterations and denote 
%     $L_{\Delta}= \frac{\max\{S_{1}, S_{2}, S_{3}, S_{4}\}}{\min\{C_{\bB},
%     C_{\bx}\}}$, $C_{\sigma}= \frac{L_{\Delta}}{L_{F}}+ 3$, and
%     $\Delta_{1}= \Psi_m^{1}- \Psi_m(\mathcal{Q}^{*})$
%     are constants independent of $T$, where \(\Psi_m^{1} = \Psi_m(\calQ^1)\).
% \begin{lemma}[Convergence rate of \Cref{alg:our_algorithm_sgdm}]
%     \label{lemma: convergence rate sgdm} Under the conditions of Lemma \ref{Lemma:
%     approximate sufficient descent property in SGDM}, for any $T \geq 1$, there exists
%     $R$ uniformly selected from $\{2,\ldots,T+1\}$ such that $(\bx^{R}, \bB
%     ^{R}, \bLmbd^{R})$ satisfies \eqref{equ:esp} with
%     \begin{align}
%         \label{Eq: the convergence rate of sgdm}\mathbb{E}[\mathrm{dist}^{2}(0, \partial \mathcal{L}(\bx^{R}, \bB^{R}, \bLmbd^{R}))] \leq \frac{L_{\Delta}\Delta_{1}}{T}+ C_{\sigma}\frac{\sigma^{2}}{b_{k}}.
%     \end{align}
% \end{lemma}
\begin{proof}
    \blue{Under \eqref{Eq: the parameters conditions of dualhash-sgdm in suppl}, $\eta_k=\eta/L_F$, $\alpha_k=\alpha$, and $\beta_k=\beta$ are constant, so $S_3^k=S_3$ and $S_4^k=S_4$. Since the four increment terms are nonnegative, \eqref{Eq: the stationarity error bound in SGDM Lemma} gives}
    {\small
    \begin{align}\label{Eq: the stationarity error bound in SGDM rearrange}
        \mathbb{E}[\mathrm{dist}^{2}(\bm{0}, \mathcal K(\bx^{k}, \bB^{k}, \bLmbd^{k}))] \leq \blue{\max\{S_{1}, S_{2}, S_{3}, S_{4}\}} \EE [\Delta_{\bB}^{k+1}+ \Delta_{\bB}^{k}+ \Delta_{\bx}^{k+1}+ \Delta_{\bx}^{k}] + \blue{\frac{3 \sigma^{2}}{b_{k}}}.\end{align}}
    Similarly, from \eqref{Eq: the descent property of Lyapunov func in SGDM}, we
    obtain
    \[
        \EE [\Psi_m(\mathcal{Q}^{k+1})] - \EE[\Psi_m(\mathcal{Q}^{k})]
        \leq -\min \{ C_{\bB}, C_{\bx}\}\EE[\Delta_{\bB}^{k+1}+ \Delta_{\bB}^{k}+
        \Delta_{\bm{x}}^{k+1}+ \Delta_{\bm{x}}^{k}] + \frac{\sigma^{2}}{b_{k}L_{F}},
    \]
    which yields that
    \[
        \EE[\Delta_{\bB}^{k+1}+ \Delta_{\bB}^{k}+ \Delta_{\bm{x}}^{k+1}+ \Delta_{\bm{x}}
        ^{k}] \leq \frac{1}{\min \{ C_{\bB}, C_{\bx}\}}\left( \EE [\Psi_m(\mathcal{Q}
        ^{k})] - \EE[\Psi_m(\mathcal{Q}^{k+1})] + \frac{\sigma^{2}}{b_{k}L_{F}}
        \right)
    \]
    Substituting this inequality into \eqref{Eq: the stationarity error bound in SGDM rearrange},
    summing over $k = 1, \cdots,T$, and dividing by $T$, we obtain
    \begin{align*}
        &\frac{1}{T}\sumT \mathbb{E}[\mathrm{dist}^{2}(\bm{0}, \mathcal K(\bx
        ^{k}, \bB^{k}, \bLmbd^{k}))]\\
        &\leq \frac{L_{\Delta}}{T}\sumT
        \left(\EE [\Psi_m(\mathcal{Q}^{k})] - \EE [\Psi_m
        (\mathcal{Q}^{k+1})] \right) + \blue{\left( 3 + \frac{L_{\Delta}}{L_F}\right) \frac{\sigma^{2}}{T}\sumT\frac{1}{b_{k}}},
    \end{align*}\blue{Since $\bx^0=\bx^1$ and $\bB^0=\bB^1$, the negative increment terms give $\Psi_m(\calQ^1)\le\calL(\bx^1,\bB^1,\bLmbd^1)$, so $\EE[\Psi_m(\calQ^1)]-\Psi_m^*\le\Delta_1$.} Denoting
    $a_{k}= \mathbb{E}[\mathrm{dist}^{2}(\bm{0}, \mathcal K(\bx^{k}, \bB
    ^{k}, \bLmbd^{k}))]$, we have
    \[
        \frac{1}{T}\sumT a_{k}\leq \frac{L_{\Delta}\Delta_{1}}{T}+ \blue{\left( 3 + \frac{L_{\Delta}}{L_F}
        \right) \frac{\sigma^{2}}{T}\sumT\frac{1}{b_{k}}}.
    \]
    \blue{By selecting $R$ uniformly from $\{1, \ldots, T\}$, independently of the iterates, and applying the law of total expectation,}
    we obtain
    \[
        \mathbb{E}[\mathrm{dist}^{2}(\bm{0}, \mathcal K(\bx^{R}, \bB^{R}
        , \bLmbd^{R}))] \blue{= \frac{1}{T}\sumT a_{k}} \leq \frac{L_{\Delta}\Delta_{1}}{T}+ \blue{\left( 3 + \frac{L_{\Delta}}{L_F}
        \right) \frac{\sigma^{2}}{T}\sumT\frac{1}{b_{k}}},
    \] and complete the proof.
\end{proof}
% From \eqref{Eq: the convergence rate of sgdm}, to achieve $\calO(T^{-1})$
% convergence rate, we need to set $b_{k}= \Theta (T)$. Next, we set
% $b_{k}= b = c_{b}T$ for all $k \geq 1$ in this way and establish the complexity
% results of \Cref{alg:our_algorithm_sgdm} with the parameter conditions \eqref{Eq: the parameters conditions of dualhash-sgdm in suppl}.
\subsection{Proof of \Cref{Theorem: the oracle complexity in SGDM}}
\label{app:stom-oracle-proof}
% \setcounter{proof}{0}
\begin{proof}
    When $b_{k}= b = c_{b}T$, \blue{we have $\frac{1}{T}\sumT\frac{1}{b_{k}}=\frac{1}{c_{b}T}$, and hence \Cref{lemma: convergence rate sgdm} gives}
    \[
        \mathbb{E}[\mathrm{dist}^{2}(\bm{0}, \mathcal K(\bx^{R}, \bB^{R}
        , \bLmbd^{R}))] \leq \frac{1}{T}\left( L_{\Delta}\Delta_{1}+ C_{\sigma}
        \frac{\sigma^{2}}{c_{b}}\right).
    \]
    With the conditions on $T$ in \eqref{Eq: the iterate number upper bound of algorithm SGDM},
    this gives
    $\mathbb{E}[\mathrm{dist}^{2}(\bm{0}, \mathcal K(\bx^{R}, \bB^{R}
    , \bLmbd^{R}))] \leq \varepsilon^{2}$. Then, the oracle complexity
    of SFO is
    \[
        \text{Complexity}= (T-1)b_{k}=\calO(T^{2}) = \calO(\varepsilon^{-4}).
    \]
\end{proof}
\section{Complexity Analysis of DualHash-StoRM (\texorpdfstring{\Cref{alg:our_algorithm_storm}}{Algorithm 2})}
\label{app:storm-complexity}
We adopt the same notation as in \Cref{app:stom-complexity}. 
\subsection{\blue{Proof of \Cref{Lemma: One-iteration progress of x in STORM}}}
\label{app:storm-primal-progress-proof}
\begin{proof}
    \blue{By the $L_F$-smoothness of $F(\cdot,\bB^k)$, the update \eqref{Eq: the update of x in STORM}, and Young's inequality,}
    \blue{\begin{align}
        &\calL(\bx^{k+1},\bB^k,\bLmbd^k)-\calL(\bx^k,\bB^k,\bLmbd^k)\notag\\
        &=F(\bx^{k+1},\bB^k)-F(\bx^k,\bB^k)\notag\\
        &\leq\eta_k\langle\be^k,\bg_s^k\rangle-\eta_k\Bigl(1-\frac{\eta_kL_F}{2}\Bigr)\|\bg_s^k\|^2\notag\\
        &\leq\frac{\eta_k}{2}\|\be^k\|^2-\frac{\eta_k}{2}(1-\eta_kL_F)\|\bg_s^k\|^2\notag\\
        &=\frac{\eta_k}{2}\|\be^k\|^2-(\eta_k^{-1}-L_F)\Delta_{\bx}^{k+1},\notag
    \end{align}}
    \blue{where we used $\nabla_{\bx}F(\bx^k,\bB^k)=\bg_s^k-\be^k$ and $2\Delta_{\bx}^{k+1}=\eta_k^2\|\bg_s^k\|^2$.}
\end{proof}

\subsection{Proof of \Cref{Lemma: approximate sufficient descent property storm}}
\label{app:storm-descent-proof}
\begin{proof}
    Substituting the one-iteration progress \eqref{Eq: the recursive relation of L on x}
    and \eqref{item: B and Lambda} into \eqref{Eq: One-iteration progress decomposition}, we obtain
    \begin{align}
             & \calL (\bx^{k+1}, \bB^{k+1}, \bLmbd^{k+1}) - \calL (\bx^{k}, \bB^{k}, \bLmbd^{k}) \nonumber                                                                                                                                          \\
        \leq & \,\frac{\eta_{k}}{2}\| \be^{k}\|^{2}- (\eta_{k}^{-1}- L_{F}- K_{3}) \Delta_{\bx}^{k+1}+ K_{1}\Delta_{\bB}^{k+2}+ K_{4}\Delta_{\bx}^{k+2}- K_{2}\Delta_{\bB}^{k+1}+ K_{3}\Delta_{\bB}^{k}.\label{Eq: One-iteration progress in storm}
    \end{align}
    Applying the definition of $\Psi_r$ from \eqref{Eq: the Lyapunov function in STORM} into \eqref{Eq: One-iteration progress in storm}, we rearrange the above inequality
    \begin{align}
        \label{the descent property of Lyapunov func in STORM first} & \Psi_r(\mathcal{Q}^{k+1}) - \Psi_r(\mathcal{Q}^{k}) \notag                                                                                         \\
        \leq                                                         & - (C_{1}- K_{1}) \Delta_{\bm{B}}^{k+2}-(K_{2}- C_{1}- C_{2}) \Delta_{\bm{B}}^{k+1}- (C_{2}- K_{3}) \Delta_{\bm{B}}^{k}- (C_{3}-K_{4}) \Delta_{\bx}^{k+2}\notag \\
                                                                     & \blue{-(\eta_{k}^{-1}- L_{F}- K_{3}- C_{3}) \Delta_{\bx}^{k+1}+ \frac{\eta_{k}}{2}\| \be^{k}\|^{2}}.
    \end{align}
    By substituting the Lyapunov function coefficients \eqref{Eq: the coefficients of the Lyapunov function in SGDM} into \eqref{the descent property of Lyapunov func in STORM first} \blue{and taking full expectation}, we obtain \eqref{Eq: the descent property of Lyapunov func in storm}. With $\tau$ satisfying \blue{$0<\tau<\frac{\sqrt{30}-2}{13L_F}$}, \Cref{lem:lyapunov-coef-positivity} guarantees that $C_{\bB}>0$.
\end{proof}

% \subsection{Proof of \Cref{Lemma: Weighted error accumulation}}
% \label{Subsubsec: proof Weighted error accumulation}
\subsection{\blue{Lower Bound of the StoRM Potential}}
\label{app:storm-lower-bound}
\begingroup\color{blue}
\begin{lemma}[Lower bound of the StoRM potential]
\label{Lemma: lower bound of W}
Under Assumptions~\ref{Ass: basic regularity} and~\ref{Ass: main assumptions} and the parameter conditions \eqref{Eq: the parameters conditions in dualhash-storm}--\eqref{Eq: the constants for storm}, for every $k\ge1$,
\[
W_k:=\EE[\Psi_r(\calQ^k)]+\frac{\eta_T}{\rho_T}\EE\|\be^k\|^2\ge\Psi_r^*:=\underline{\ell}-c_1',
\]
where $c_1'$ is defined in \eqref{eq:lb-common} and is independent of $T$.
\end{lemma}

\begin{proof}
By \eqref{Eq: the update of x in STORM} and \eqref{eq:lb-gradient},
\[
C_3\Delta_{\bx}^{k+1}
=\frac{\tau L_F^2\eta_T^2}{4}\|\nabla_{\bx}F(\bx^k,\bB^k)+\be^k\|^2
\le\tau L_F^3\eta_T^2\bigl(F(\bx^k,\bB^k)-\underline{\ell}\bigr)+\frac{\tau L_F^2\eta_T^2}{2}\|\be^k\|^2 .
\]
Since $\tilde L_F\ge L_F$, $\eta_T\tilde L_F\le\frac12$ and $\rho_T\le\frac12$ give $\eta_T\le\frac{1}{2L_F}$ and $\eta_T\rho_T\le\frac{1}{4L_F}$. Therefore
\[
\tau L_F^3\eta_T^2\le\frac t4<1,
\qquad
\frac{\tau L_F^2\eta_T^2}{2}=\frac{\tau L_F^2\eta_T\rho_T}{2}\cdot\frac{\eta_T}{\rho_T}\le\frac t8\cdot\frac{\eta_T}{\rho_T}<\frac{\eta_T}{\rho_T}.
\]
Using \eqref{eq:lb-common} and $C_2\Delta_{\bB}^k\ge0$, we obtain
\[
W_k\ge\underline{\ell}+\left(1-\frac t4\right)\EE\bigl[F(\bx^k,\bB^k)-\underline{\ell}\bigr]+\left(1-\frac t8\right)\frac{\eta_T}{\rho_T}\EE\|\be^k\|^2-c_1'\ge\underline{\ell}-c_1' .
\]
\end{proof}
\endgroup

\subsection{\blue{Auxiliary Results for DualHash-StoRM}}
\label{app:storm-stationarity-proof}
We provide a recursive bound on the gradient error vector $\{\be^{k}\}$ with analysis comparable to those in \cite[Lemma~2\blue{.2}]{xu2023momentum} and \cite[Lemma 8]{shi2025momentum}. \blue{We provide the proof for completeness.}
\begin{lemma}[Recursive bound on gradient error]
    \label{Lemma: Recursive bound on gradient error} Under \Cref{Ass: main assumptions}\blue{(ii) and \Cref{Ass: the mean squared smoothness of F}},
    for any $k \geq 1$, it holds that
    \begin{align}
        \blue{\EE}[\| \be^{k+1}\|^{2}] \leq \left(1 - \rho_{k}\right)^{2}\blue{\EE}[\| \be^{k}\|^{2}] + \frac{2 \rho_{k}^{2}\sigma^{2}}{\blue{b_{k+1}}}+ \frac{4(1 - \rho_{k})^{2}L_{F}^{2}}{\blue{b_{k+1}}}\blue{\EE}\left[\Delta_{\bx}^{k+1}\blue{{}+\Delta_{\bB}^{k+1}}\right]. \notag
    \end{align}
\end{lemma}
\begin{proof}
    \blue{By \eqref{Eq: the stochastic approximate in STORM}, the definition of $\be^{k}$, and $\EE_{k+1}[\tilde{\nabla}_{\bx}F(\bx^{j},\bB^{j};\calJ_{k+1})]=\nabla_{\bx}F(\bx^{j},\bB^{j})$ for $j\in\{k,k+1\}$, we obtain
    \begin{align}
        \EE_{k+1}[\|\be^{k+1}\|^{2}]
        \leq & \,(1-\rho_k)^{2}\|\be^{k}\|^{2}+\frac{2\rho_k^{2}\sigma^{2}}{b_{k+1}}\notag\\
        & \,+\frac{2(1-\rho_k)^{2}}{b_{k+1}}\cdot\frac1n\sum_{i=1}^{n}\|\nabla_{\bx}F_i(\bx^{k+1},\bb_i^{k+1})-\nabla_{\bx}F_i(\bx^{k},\bb_i^{k})\|^{2}\notag\\
        \leq & \,(1-\rho_k)^{2}\|\be^{k}\|^{2}+\frac{2\rho_k^{2}\sigma^{2}}{b_{k+1}}+\frac{4(1-\rho_k)^{2}L_F^{2}}{b_{k+1}}\left(\Delta_{\bx}^{k+1}+\Delta_{\bB}^{k+1}\right)\quad\text{a.s.},\notag
    \end{align}
    where the first inequality follows from $\|\bm a+\bm b\|^{2}\leq2\|\bm a\|^{2}+2\|\bm b\|^{2}$, \Cref{Ass: main assumptions}(ii), and the mini-batch variance bound under uniform sampling without replacement, and the last inequality follows from \Cref{Ass: the mean squared smoothness of F}. Taking the full expectation completes the proof.}
\end{proof}

We now establish a weighted error accumulation bound for the variance-reduced gradient estimator in \Cref{alg:our_algorithm_storm}.
\begin{lemma}[Weighted error accumulation bound]
    \label{Lemma: Weighted error accumulation} Under Assumption\blue{s}~\ref{Ass: basic regularity}, \ref{Ass: main assumptions}, \blue{and~\ref{Ass: the mean squared smoothness of F}} and parameter conditions \blue{\eqref{Eq: the parameters conditions in dualhash-storm}--\eqref{Eq: the constants for storm}}, for any $T \geq 1$, the following weighted cumulative error bound holds:
    \begin{align}
        \blue{\frac{1}{T}\sumT\EE\Bigl[C_{\bB}\bigl(\Delta_{\bB}^{k+1}+\Delta_{\bB}^{k}\bigr)+\frac{1}{4\eta_T}\Delta_{\bx}^{k+1}+\frac{\eta_T}{2}\|\be^{k}\|^{2}\Bigr]\leq\frac{1}{T}\Bigl(\tilde{\Delta}_{1}+\frac{\eta_T\sigma^{2}}{\rho_Tb_1}\Bigr)+\frac{2\eta_T\rho_T\sigma^{2}}{b},}
        \label{Eq: weighted error accumulation bound}
    \end{align}
    where $\tilde{\Delta}_{1}=\blue{\calL(\bx^{1},\bB^{1},\bLmbd^{1})}-\Psi_r^{*}$.
\end{lemma}
\begin{proof}
    Under \blue{\eqref{Eq: the parameters conditions in dualhash-storm}--\eqref{Eq: the constants for storm}, $\eta_T\tilde L_F\le\frac12$, so} we have a lower bound for $C_{\bx}^{k}$ with $C_{\bx}^{k}\geq\frac{1}{2\blue{\eta_T}}>0$. Substituting this inequality into \eqref{Eq: the descent property of Lyapunov func in storm} yields
    \begin{align}
        \EE[\Psi_r(\calQ^{k+1})]-\EE[\Psi_r(\calQ^{k})]\leq\blue{-C_{\bB}\EE[\Delta_{\bB}^{k+1}+\Delta_{\bB}^{k}]}-\frac{1}{2\blue{\eta_T}}\EE[\Delta_{\bx}^{k+1}]+\frac{\blue{\eta_T}}{2}\EE[\|\be^{k}\|^{2}].
        \label{Eq: descent of Psi_r in weighted error bound}
    \end{align}
    \blue{To handle the term $\Delta_{\bB}^{k+1}$ in \Cref{Lemma: Recursive bound on gradient error}, let $r=\tau\gamma/n$ and
    \[
        H_k=(1-r)\|\tau(\bLmbd^{k}-\bLmbd^{k-1})-(\bB^{k}-\bB^{k-1})\|^{2}+r\tau^{2}\|\bLmbd^{k}-\bLmbd^{k-1}\|^{2}
    \]
    with $\bLmbd^{0}=\bLmbd^{1}$. Since the initialization sets $\bB^{0}=\bB^{1}$ and $\bLmbd^{0}=\bLmbd^{1}$, we have $H_1=0$ and $H_k\geq0$. By the $\bB$- and $\bLmbd$-updates, we obtain}
    \begin{align}
        \blue{\Delta_{\bB}^{k+1}\leq\frac{12}{r}\bigl(H_k-H_{k+1}\bigr)+G\,\Delta_{\bx}^{k+1},}
        \label{Eq: tracking of B increments}
    \end{align}
    \blue{where the inequality follows from $r\le\tau L_F<\frac{\sqrt{30}-2}{13}<\frac27$ and the entrywise monotonicity and $1$-Lipschitz continuity of $\operatorname{prox}_{\tau^{-1}h^{*}}$, which give $H_{k+1}-H_k\leq-\frac{r}{8}\bigl(\|\bB^{k}-\bB^{k-1}\|^{2}+\|\tau(\bLmbd^{k}-\bLmbd^{k-1})-(\bB^{k}-\bB^{k-1})\|^{2}\bigr)+146rG_{\mathcal{D}}^{2}\Delta_{\bx}^{k+1}$.}
    From \Cref{Lemma: Recursive bound on gradient error} \blue{and \eqref{Eq: tracking of B increments}}, we obtain
    \begin{align}
        &\blue{\EE[\|\be^{k+1}\|^{2}]+\frac{48L_F^{2}}{br}\EE[H_{k+1}]}\notag\\
        &\blue{\leq(1-\rho_T)\EE[\|\be^{k}\|^{2}]+\frac{48L_F^{2}}{br}\EE[H_{k}]+\frac{2\rho_T^{2}\sigma^{2}}{b}+\frac{4L_F^{2}(1+G)}{b}\EE[\Delta_{\bx}^{k+1}],}
        \label{Eq: error recursion with tracking}
    \end{align}
    where the inequality follows from \blue{$(1-\rho_T)^{2}\leq1-\rho_T$}. \blue{Multiplying \eqref{Eq: error recursion with tracking} by $\frac{\eta_T}{\rho_T}$, adding it to \eqref{Eq: descent of Psi_r in weighted error bound}, and applying $\frac{\eta_T}{\rho_T}\cdot\frac{4L_F^{2}(1+G)}{b}\leq\frac{1}{4\eta_T}$, which is equivalent to the lower bound on $\rho$ in \eqref{Eq: the constants for storm},} yields the following expected descent property of the Lyapunov function $\Psi_r$:
    \begin{align}
        &\blue{\EE\Bigl[\Psi_r(\calQ^{k+1})+\frac{\eta_T}{\rho_T}\Bigl(\|\be^{k+1}\|^{2}+\frac{48L_F^{2}}{br}H_{k+1}\Bigr)\Bigr]-\EE\Bigl[\Psi_r(\calQ^{k})+\frac{\eta_T}{\rho_T}\Bigl(\|\be^{k}\|^{2}+\frac{48L_F^{2}}{br}H_{k}\Bigr)\Bigr]}\notag\\
        &\blue{\leq-\EE\Bigl[C_{\bB}\bigl(\Delta_{\bB}^{k+1}+\Delta_{\bB}^{k}\bigr)+\frac{1}{4\eta_T}\Delta_{\bx}^{k+1}+\frac{\eta_T}{2}\|\be^{k}\|^{2}\Bigr]+\frac{2\eta_T\rho_T\sigma^{2}}{b}.}
        \label{Eq: expected descent of Psi_r with error}
    \end{align}
    Finally, summing up \eqref{Eq: expected descent of Psi_r with error} from $k=1$ to $T$ and dividing by $T$, \blue{together with $H_1=\Delta_{\bB}^{1}=0$, $\EE[\|\be^{1}\|^{2}]\leq\sigma^{2}/b_1$, $H_{T+1}\geq0$, and \Cref{Lemma: lower bound of W},} yields \eqref{Eq: weighted error accumulation bound} and completes the proof.
\end{proof}

We derive an upper stationarity error bound for the iterates generated by \Cref{alg:our_algorithm_storm}.
\begin{lemma}[Stationarity error bound]
    \label{Lemma: Stationarity error bound} Under \Cref{Ass: basic regularity}(ii),
    for any $k \geq 1$, it holds that
    \begin{align}
        \mathrm{dist}^{2}(\bm{0}, \mathcal K(\bx^{k}, \bB^{k}, \bm{\Lambda}^{k})) \leq S_{1}\Delta_{\bB}^{k+1}+ S_{2}\Delta_{\bB}^{k}+ S_{3}^{k}\Delta_{\bm{x}}^{k+1}+ 2 \| \be^{k}\|^{2}, \label{Eq: the stationarity error bound in Lemma in storm}
    \end{align}
    where the coefficients are defined as \eqref{Eq: the coefficients of stationarity error bound in SGDM}.
    % \begin{align}
    % \label{Eq: the  coefficients of stationarity error bound in storm}
    % \begin{cases}
    %     S_1 &= 4\tau^{-2} + 12,\\
    %     S_2 &=  4(1 + 3 (1 + \tau L_F )^2),\\
    %     S_3^k &= 4(\eta_k^{-2} + L_F^2+ 3 (\tau L_F)^2).
    % \end{cases}
    % \end{align}
\end{lemma}
\begin{proof}
    This proof follows a similar approach to that of \Cref{Lemma: Stationarity error bound in SGDM} 
in \Cref{app:stom-complexity}. The key difference lies in 
the estimation of $\| \bd_{\bx}^{k}\|^{2}$. To estimate $\| \bd_{\bx}^{k}\|^{2}$, we use 
the optimality condition of $\bx^{k+1}$ from \eqref{Eq: the update of x in STORM} to obtain
    \begin{align}
        \|\bd_{\bx}^{k}\|^{2}= & \,\eta_{k}^{-2}\left\| \eta_{k}\nabla_{\bx}F(\bx^{k}, \bm{B}^{k}) + 0\right\|^{2}\notag                                       \\
        =                      & \, \eta_{k}^{-2}\left\|\eta_{k}\left( \nabla_{\bx}F(\bx^{k}, \bm{B}^{k}) - \blue{\bg_s^{k}}\right) + \eta_{k}\blue{\bg_s^{k}}\right\|^{2}\notag \\
        =                      & \, \eta_{k}^{-2}\left\|\blue{-}\eta_{k}\be^{k}+ (\bx^{k}- \bx^{k+1})\right\|^{2}\notag                                                \\
        \leq                   & \, 4\eta_{k}^{-2}\Delta_{\bx}^{k+1}+ 2 \| \be^{k}\|^{2}.\label{Eq: dx in storm}
    \end{align}
    Then, substituting \eqref{Eq: dx in storm}, \eqref{Eq: dB}, and \eqref{Eq: dLambda}
    into \eqref{Eq: the stationarity error bound in proof} yields \eqref{Eq: the stationarity error bound in Lemma in storm}
    and completes the proof. 
\end{proof}
% \begin{align}
%     \begin{cases}
%         &\bd_{\bx}^k =  \nabla_{\bx} F(\bx^k, \bB^k),\\
%         &\bd_{\bB}^k = \nabla_{\bB} F(\bx^k, \bB^k) + \bm{\Lambda}^k, \\
%         &\bd_{\bm{\Lambda}}^k = \bB^k - \calG_{\bm{\Lambda}}^k.
%     \end{cases}
% \end{align} where \(\calG_{\bm{\Lambda}}^k \in \partial h^*(\bm{\Lambda}^k)\).
% Therefore,
% \begin{align}
% \label{Eq: the stationarity error bound in proof}
% \text{dist}^2(\bm{0}, \partial \calL(\bx^k, \bB^k, \bm{\Lambda}^k)) \leq \|\bd_{\bx}^k \|^2 + \| \bd_{\bB}^k\|^2 + \| \bd_{\bm{\Lambda}}^k\|^2.
% \end{align}
%     For estimating  \(\bd^k_{\bB}\), using the optimality condition of \(\bB^k\) and the smoothness of \(F\), we obtain
% \begin{align}
%     \|\bd_{\bB}^k\|^2&= \| \nabla_{\bB} F(\bx^{k+1}, \bB^k)+  \bm{\Lambda}^k + \nabla_{\bB} F(\bx^k, \bB^k) - \nabla_{\bB} F(\bx^{k+1}, \bB^k) \|^2 \notag \\
%                     &= 4 \tau^{-2} \Delta^{k+1}_{\bB}  + 4L_F^2\Delta_{\bx}^{k+1},
% \end{align} where the last inequality follows from the fact \((a + b)^2 \leq 2(a^2 + b^2)\).

% For estimating \( \|\bd^k_{\bm{\Lambda}}\|^2\),  using the optimality condition of \( \bm{\Lambda}^k \) in \eqref{Eq: the optimality condition of lbda},  we obtain
% \begin{align*}
%         \|\bd^k_{\bm{\Lambda}}\|^2 &\leq 2\tau^{2} \|\bm{\Lambda}^k - \bm{\Lambda}^{k-1} \|^2+ 2\|\bB^{k} - \bB^{k-1}\|^2.
%     \end{align*} Then, substituting \eqref{Eq: dual boundedness} in \cref{Lemma: dual boundedness} into the above inequality yield that
% \begin{align}
% \|\bd^k_{\bm{\Lambda}} \|^2 \leq 12 \Delta_{\bB}^{k+1} + 4(1 + 3 (1 + \tau L_F )^2)\Delta_{\bB}^k + 12(\tau L_F)^2\Delta_{\bx}^{k+1} .
% \end{align}

\subsection{Proof of \Cref{Lemma: convergence rate of storm}}
\label{app:storm-rate-proof}
\begin{proof}
    By taking the full expectation on both sides of \blue{\eqref{Eq: the stationarity error bound in Lemma in storm} in \Cref{Lemma: Stationarity error bound}} and averaging over $T$, we obtain
    \[
        \blue{\frac1T\sumT\EE[\mathrm{dist}^{2}(\bm{0},\mathcal K(\bx^k,\bB^k,\bLmbd^k))]
        \leq\frac1T\sumT\EE\bigl[S_1\Delta_{\bB}^{k+1}+S_2\Delta_{\bB}^{k}+S_3^k\Delta_{\bx}^{k+1}+2\|\be^{k}\|^{2}\bigr].}
    \]
    \blue{By the definition of $\tilde{L}_{\Delta}$, we have $S_1,S_2\leq\tilde{L}_{\Delta}T^{1/3}C_{\bB}$, $4\eta_TS_3^k\leq\tilde{L}_{\Delta}T^{1/3}$, and $\frac{4}{\eta_T}\leq\tilde{L}_{\Delta}T^{1/3}$.} Then, substituting \blue{\eqref{Eq: weighted error accumulation bound} from \Cref{Lemma: Weighted error accumulation}} into the above inequality yields
    \[
        \blue{\frac1T\sumT\EE[\mathrm{dist}^{2}(\bm{0},\mathcal K(\bx^k,\bB^k,\bLmbd^k))]
        \leq\frac{\tilde{L}_{\Delta}}{T^{2/3}}\Bigl(\tilde\Delta_1+\frac{\eta_T\sigma^2}{\rho_Tb_1}+\frac{2T\eta_T\rho_T\sigma^2}{b}\Bigr).}
    \]
    \blue{Substituting $\frac{\eta_T}{\rho_T}=\frac{T^{1/3}}{8\rho\eta\tilde L_F}$, $2T\eta_T\rho_T=\frac{16\rho\eta^3}{\tilde L_F}$, and $\tilde{C}_{\sigma}=\frac{\eta\tilde{L}_{\Delta}}{2\tilde L_F}$,} we obtain \eqref{Eq: the convergence rate of storm} by selecting $R$ uniformly from \blue{$\{1,\ldots,T\}$}.
\end{proof}
With this parameter choice, we proceed to analyze the oracle complexity of \Cref{alg:our_algorithm_storm}.
\subsection{Proof of \Cref{Theorem: the oracle complexity in STORM}}
\label{app:storm-oracle-proof}
\begin{proof}
    When $b_1=c_bT^{1/3}$, we have
    \[
        \EE[\mathrm{dist}^{2}(\bm{0}, \mathcal K(\bx^{R}, \bB^{R}, \bm{\Lambda}^{R}))]\le\frac{1}{T^{2/3}}\Bigl(\tilde{L}_{\Delta}\tilde{\Delta}_{1}+\tilde{C}_{\sigma}\Bigl(\frac{\sigma^{2}}{4\rho\eta^{2}c_b}+\frac{32\rho\eta^{2}\sigma^{2}}{b}\Bigr)\Bigr).
    \]
    With the conditions on $T$ in \eqref{Eq: the iterate number upper bound of algorithm VR}, this gives $\EE[\mathrm{dist}^{2}(\bm{0}, \mathcal K(\bx^{R}, \bB^{R}, \bm{\Lambda}^{R}))]\le\varepsilon^2$. Then, the oracle complexity is
    \[
        \text{Complexity}= b_{1}+ (T-1)b = \calO(T^{\frac{1}{3}}) + \calO(T) = \calO(T) = \calO(\varepsilon^{-3}),
    \]
    which completes the proof.
\end{proof}
% \[
% \eta_k = \frac{\eta}{L_F} \frac{1}{T^l}, \quad (L_F + K_3) \eta_k \leq \frac{1}{2}
% \] where \(\rho\) and \(\eta\) are independent of \(T\).
% Let us \(l = \frac{1}{3}\) and \(q = \frac{2}{3}\),
% we want \(\frac{1}{T}\sum_{k=1}^T\text{dist}^2(\bm{0}, \partial \calL^k) = \mathcal{O}(T^{-\frac{2}{3}})\)
% We denoteare constant independent of \(T\) and \(\Delta_1 = \Psi_r(\calQ^1) - \Psi_r^* \) and the definitions of the \(\Psi_r\) and \(\calQ\) are detailed in Appendix \ref{app:storm-complexity}.

% Let us discuss how to select appropriate parameters.  To ensure the parameter setting \eqref{Eq: the parameters conditions in dualhash-storm}, we set the parameters as follows
% \[
% \eta_k = \frac{\eta}{T^l\tilde{L}_{F}},\quad \rho_k = \frac{8\rho\eta^2}{T^q}
% \] where \(q \leq 2l\), \(0< \eta \leq \frac{1}{2}\) and \(1 \leq \rho \leq \frac{1}{8\eta^2}\) are given constants that are independent of \(T\).
% Then, the upper bound of \eqref{Eq: the convergence rate of algorithm VR} is
% \[
% \text{Rate} = \calO (\max \{ T^{l-1}, T^{q-p-1}, T^{-2q}\})
% \]
% We can choose \(q = 2l = \frac{2}{3}\) and \(p=\frac{1}{3}\), then the upper bound is of the order \(\calO(T^{-\frac{2}{3})}\)
% \subsection{Discussion}
